\chapter{}

Então, ele finalmente desabou.
Ao longo desses meses passara pela morte da mãe, Dona Yolanda e pela de um meio-irmão, Mário.
Todo nosso dinheiro se fora e ele passava os dias numa atitude que eu chamava de ``urubu molhado'', cabeça metida entre os ombros, calado, vagando pela casa embrulhado num casaco de lã, como um velho, mergulhado na mais profunda crise de autocomiseração.
Dia após dia, aquela visão foi me dando nos nervos e fez subir à tona meus piores recalques.
Tenho enorme dificuldade em entender e lidar com estados prolongados de infelicidade.
Compaixão, para mim, é remédio de uso tópico e controlado, só aplicável aos casos em que se reconhece, naquele a quem o sentimento se dirige, alguém menos capaz de assumir responsabilidade pela própria vida.
Mesmo com tudo que lhe sucedera de ruim, eu não queria enxergar no Paulo alguém nessas condições.

Nessa altura, a nossa situação era bastante complicada: nossos filhos iam terminando seus estudos ou já trabalhavam, namoravam e pensavam em casar.
Estávamos vivendo numa cidade provincianamente preconceituosa que eu odiava e cercados por uma família que tinha uma enorme dificuldade em entender nosso modo de vida, o que me mantinha em permanente defensiva para evitar que críticas e intromissões indevidas viessem a atingir Paulo.
Meus irmãos até me ajudariam, mas eu os conhecia o suficiente para saber que essa ajuda seria muito condicional.
Portanto, era uma alternativa fora de questão.
Eu tinha certeza que ali era o pior lugar para Paulo estar naquele momento.
E, por outro lado, tinha a esperança de que nos lugares por onde tínhamos andado, nos quais Paulo sempre encontrara aceitação, ele pudesse descobrir uma oportunidade de recomeçar e recuperar-se daquela sucessão de revezes.
Comecei cuidadosamente a sugerir-lhe a ideia.

Se há uma coisa que em Paulo é aguçada é a capacidade de perceber intenções ocultas e manipulação.
E ele detectou isso imediatamente nas minhas palavras.
Havia alguma coisa que eu não queria dizer.
E, de fato havia, embora naquele instante ainda não soubesse com clareza o quê.
Eu acreditava sinceramente que ir embora para o norte era o melhor que ele tinha a fazer.
Eu conhecia a fundo o mal que o ambiente araraquarense podia causar a alguém nas condições dele.
O que não disse é que, chegada a hora de os nossos filhos começarem sua própria caminhada na vida, eu faria das tripas coração para que nada anuviasse esse momento, para que nada comprometesse a alegria e a confiança que deviam animar seus primeiros passos nessa trilha.
 Marcelo, Fernando e Paula haviam de se casar um após o outro ao longo de um ano.
No estado de espírito em que Paulo se encontrava, eu tinha certeza, seria impossível fazer o que eu sentia que precisava fazer por eles.

Ele acabou indo embora magoado, sentindo-se rejeitado e eu sabia que esse era o seu pior trauma - o da rejeição.
Mas, eu atravessaria uma fase, daí por diante, durante a qual teria que esquecer até de mim mesma, se quisesse encarar o que tinha pela frente.
Se um dia ele pudesse entender isso, ótimo, senão, paciência, não havia o que fazer, pensei.
Antes de partir, Paulo pediu a Tomaz e Marcelo que assumissem o sustento da casa, até que ele pudesse se aprumar.
Marcelo fora contratado pela EMBRAER e, junto com a noiva, que trabalhava como farmacêutica, empenhava-se em tornar possível o casamento de ambos.
Fernando se formara e, após trabalhar um ano no Brasil para juntar dinheiro, viajara com uma bolsa de estudos para a França e lá travava sua própria luta para sobreviver.
Paula diplomara-se em Pedagogia e, para nossa alegria, passara num concurso para ser promotora cultural na unidade do SESC recém-construída em Araraquara.
Vicente terminava seu curso de Engenharia e trabalhava em estágios sucessivos, desde o primeiro ano da faculdade, para se manter.
Tomaz, já empregado como mestre cervejeiro na Antártica e ainda solteiro, achou que era o que estava em melhor condição de assumir a responsabilidade de nos amparar.
Usou, inclusive, suas economias para associar-se ao pai em Santarém, num negócio de exportação de colágeno extraído da bexiga de peixes.

Eu tinha passado por uma sucessão de problemas de saúde que começaram ainda quando eu cuidava da fazenda de peixes: quebrei o outro pé, contraí toxoplasmose, operei o joelho direito e enfim acabei em casa, reduzida às panelas e ao tanque.
Dali me tirou João que, em vias de morar uma temporada nos Estados Unidos, pediu que eu o ajudasse no atendimento do consultório até a sua partida.
Nenhum de nós dois estava muito bem, psicologicamente falando, ambos cheios de problemas pessoais e ainda desgastados pelas brigas de família, no decorrer das quais tínhamos assumindo a liderança da dissidência.
Mas, eu não estava em condições de abrir mão de qualquer ajuda.
Fui trabalhar de atendente do João e descobri que a maior vantagem de nascer Filpi é não correr o risco de casar com um deles e que, ao menos para mim, ter sido educadora foi muito melhor do que ter optado pela Medicina.

Dizem, e eu concordo plenamente, que só vale a pena fazer na vida aquilo que exija irmos além de nós mesmos.
Não obstante todas as dificuldades, nós conseguimos vencer o desafio de viver três noivados, três casamentos, a formatura do Vicente com toda a alegria que esses acontecimentos mereciam, mesmo porque, por mérito deles, cada evento realizou-se de acordo com o desejo de cada um e o principal, no que diz respeito aos noivados e casamentos, as escolhas não podiam ter sido melhores e foram a grata recompensa para o esforço feito.
Pouco antes da formatura, Vicente teve a perna gravemente fraturada num jogo de futebol e esse foi um novo e inesperado desafio a vencer, para nós dois.
Mas, no dia de receber o diploma, a muleta fez sucesso e ele acabou sendo paparicado pelos colegas e até pela organização do evento.
O incidente serviu também para demonstrar a capacidade de conquistar amigos do nosso caçula.

Tivemos, depois disso, outros problemas sérios a resolver.


Até hoje não sei até que ponto meu empenho em ocultar dos meninos os piores aspectos das nossas dificuldades foi positivo para eles.
Na realidade, essas coisas não são tão fáceis de esconder, porém, eu acreditava que enfrentando tudo com a discrição e naturalidade possíveis, eu evitaria angustiá-los com questões cuja solução estava fora do alcance deles buscar.
A meu ver, já bastavam as tensões próprias de um início de vida que eles estavam vivendo.
No entanto, é possível que eu tivesse, afinal, com a minha atitude, gerado mais insegurança do que tranquilidade para eles.
Isso ficou claro quando Tomaz, Marcelo e eu tivemos que convencer Vicente a desistir do seu projeto de esperar por um hipotético estágio na Rede Ferroviária Federal e buscar um trabalho melhor remunerado e imediato.
Ao ser inteirado dos fatos que motivavam aquela conversa, ele reagiu negativamente e manifestou revolta por ter sido mantido na ignorância da nossa real situação.
Apesar disso, acatou nosso pedido.
Mas custou a conformar-se com o modo como as coisas aconteceram e acho que isso teve a ver com certa instabilidade que o levou a trocar sucessivamente de empregos no período que se seguiu.

Enquanto isso, lá no norte, Paulo dava a impressão de não estar nada bem.
Ele vencera os desafios técnicos de produzir colágeno a partir de bexiga de peixe e o negócio, mais uma vez, foi exemplarmente montado e organizado.
Mas o mercado era difícil, dominado pela concorrência desleal dos chineses e os parceiros ingleses eram terrivelmente lentos nas decisões.
Tomaz já esgotara suas economias e, compreensivelmente, não estava disposto a recorrer a bancos.
Vicente, entre uma mudança e outra de emprego, foi lá ajudar o Paulo, a pedido dos irmãos mais velhos.
Ficou até onde aguentou, mas a experiência foi duríssima.
A partir daí, acho que amadureceu e cresceu.

Paula foi a última a casar.
Ela e o marido moraram uns tempos comigo, até decidirem o que fariam.
Os planos eram de residir na Alemanha por seis meses, mas Pedro foi brilhantemente aprovado num concurso para reger a cadeira de Termodinâmica na Faculdade de Engenharia Química da Universidade de São Paulo e eles acabaram se mudando para a capital.

Um dia, eu me vi sozinha na casa.
Não havia mais ninguém.
Sentada à beira de uma cisterna abandonada, no quintal, eu me perguntei se aquilo era tudo.
Se o que me restava era ficar ali, naquela cidade que eu detestava, tomando conta daqueles moveis, daquelas louças, daquelas lembranças de um tempo que tinha acabado.
Será que havia alguma coisa na vida que eu ainda pudesse fazer?

Tia Glória, que perto da morte se tornara sábia, pareceu entender melhor que qualquer pessoa o que se passava dentro de mim.
Telefonou um dia, logo pela manhã, contando que vira na televisão um aviso de que haveria um concurso para professora de História na Escola Industrial.
Armei-me de coragem e fui lá me inscrever.

A fila de candidatos era enorme e entre eles reconheci alguns ex-alunos meus.
Atrás de mim, um rapaz e uma moça reclamavam da coincidência de datas entre o concurso e a entrega dos seus trabalhos finais na Pós-Graduação da UNESP.
Fiquei intimidada.
Que chance eu havia de ter, depois de tantos anos afastada das salas de aula? Que diriam meus ex-alunos, se eu fosse reprovada? Decidi ir em frente, mesmo assim.
Mais por baixo do que eu estava me sentindo não haveria de ficar.
Inscrevi-me, anotei os temas selecionados para as provas escritas e para a aula diante da banca examinadora e fui para casa estudar.

Um mês depois fiz as provas escritas e, em seguida, apresentei-me para a aula.
O tema tinha sido sorteado na véspera.
Fui das últimas a comparecer diante da banca.
Seria uma demonstração de quinze minutos apenas.
Comecei e, então, não pude deixar de pensar na semelhança que há entre um professor e um ator.
De repente, foi como se eu tivesse diante de mim não um grupo de avaliadores, mas alunos.
Trinta ou quarenta minutos mais tarde, lembrei-me de parar e perguntar se meu tempo não tinha se esgotado.
Respondeu-me o presidente da banca:

\textit{``-- Por mim, pode continuar até quando quiser.
Estou gostando muito.
Fazia bastante tempo que eu não assistia a uma aula dada por um professor de verdade.''}

Saí de lá feliz.
Afinal, eu ainda vivia! Semanas depois, Paulo, que passava uns dias em casa, chegou da rua anunciando que acabara de ler a publicação do resultado: eu fora a primeira colocada no concurso e só precisava agora aguardar a nomeação.

Enquanto isso, meu antigo diretor, Padre Darly, sabendo que os filhos já não moravam mais comigo, assediava-me insistentemente para retornar a Teresina e ao meu trabalho de Coordenadora no Colégio Diocesano.
Um novo diretor tinha sido nomeado para lá e precisava de alguém experiente para assessorá-lo.
Pedi-lhe um tempo para pensar.

Minha mãe, além da Tia Glória, era quem mais tinha ficado feliz com minha aprovação no concurso.
Para ela, era a certeza de que eu permaneceria para sempre em Araraquara.
Depois de tudo que acontecera no decorrer da partilha, restara um estranhamento entre nós, os irmãos, e ela sofria com isso.
Sobretudo porque, incompreensivelmente, o filho que ela mais favorecera era o que menos a procurava.
João e Paulo também se retraíram.
E Maria Lúcia, sua outra permanente preocupação, andava muito revoltada.
Reginaldo achara um jeito de tirá-la também da sociedade na fábrica e, ainda por cima, deixando-a com a certeza de ter sido muito prejudicada no acerto de contas.
Maria Lúcia estava convencida de que mamãe, mais uma vez, calara-se para não afrontar uma decisão do filho querido e não perdia uma oportunidade de acusá-la por esta omissão.
Lectícia apegou-se, então, a mim.
Nossa casa ficava a duas quadras da dela e a ausência de Paulo animava-a a estar comigo a toda hora.


Por outro lado, Tia Glória também precisava de mim.
Com pena da minha prima Laura que, a poder de ansiolíticos, dividia seu tempo entre um marido complicado, dois filhos problemáticos, um pai genioso e a tia inválida, tomei por missão ajudá-la a cuidar desta última.
Mas Tia Glória não viveu muito.
Logo se foi, mansa, quieta e dignamente como vivera seus últimos anos.
Portanto, só uma coisa me prendia a Araraquara -- minha mãe.

Com meu amigo Pires eu aprendera muito sobre a natureza da duvidosa generosidade que leva uma pessoa a esquecer reiteradamente os seus interesses em favor do bem-estar dos outros.
Tinha aprendido a desmistificar aquele intrincado sistema de compensações cujo objetivo único, na verdade, é confirmar os envolvidos nas suas posições equivocadas, evitando o risco e o incômodo de tentar uma mudança.
Lamentei por minha mãe, por Paulo, por meus irmãos, mas eu tinha que fazer a travessia.
Eu tinha que escolher continuar a viver.
Foi o que fiz, com plena consciência e sem nenhuma culpa, confesso.
Abri mão da vaga conquistada na Escola Industrial, aceitei o convite de Padre Darly e voltei para Teresina sozinha, deixando atrás de mim um rastro de perplexidade.
Pela primeira vez, em toda minha vida, eu seria Teresa.
Nem Filpi, nem Mesquita Sampaio.
Teresa, só.
Por minha conta e risco.
Precisava saber de fato quem era eu, quanto podia, quanto valia.
